%% ============================================================
%%  Section 0: General Response & Summary of Revisions
%% ============================================================
\reviewersec{General Response \& Summary of Major Revisions}

Dear Editor-in-Chief, Associate Editor, and Reviewers,

We express our sincere gratitude to the Editorial Office and the Reviewers for their meticulous, constructive, and insightful evaluations of our manuscript (\textbf{BMEF-D-26-00169}). The feedback has highlighted crucial areas for methodological rigor, statistical discipline, and clinical precision. 

In strict adherence to the reviewers' comments, we have performed comprehensive additional experiments, conducted extensive ablations, rigorously re-evaluated our statistical framework, and substantially revised the manuscript. 

Below is an executive summary of the major revisions implemented:

\begin{enumerate}[leftmargin=1.5em, itemsep=4pt]
  \item \textbf{Circularity Elimination \& Prompt Conditioning Verification (R3-1, R2-5, R2-16):}
  We have fully disclosed the exact prompt template and all clinical features (MMSE, CDR, FAQ, GDS, NPI-Q, age, sex, weight). To definitively eliminate concerns regarding label circularity, we trained and benchmarked two dedicated clinical classification baselines (\emph{Clinical-only MLP} and \emph{MRI + Clinical MLP}), conducted complete prompt-shuffling/mismatch experiments, and performed tracer-only synthesis without clinical conditioning. The results confirm that DIReCT++ synthesizes true anatomy-grounded PET pathology rather than encoding diagnostic metadata.

  \item \textbf{Data Imputation Transparency \& Table S1 Corrections (R3-2):}
  We thoroughly re-examined Table~S1, reported exact per-variable missingness across all diagnostic cohorts, documented the median-imputation protocol within training folds, and clarified the demographic distributions to ensure total transparency.

  \item \textbf{Prognosis vs. Cross-Sectional Stratification Reframing (R3-3, R2-4, R2-20):}
  Acknowledging the cross-sectional nature of the primary ADNI dataset and the absence of longitudinal time-to-conversion endpoints, we have conscientiously removed overreaching claims of ``prognostic prediction'' and ``prognostic stratification'' across the Title, Abstract, Introduction, and Discussion, focusing precisely on cross-sectional severity differentiation (EMCI vs.\ LMCI).

  \item \textbf{Subject-Level Equivalence \& Amyloid Status Agreement (R3-4, R2-8, R2-10):}
  We have abandoned ungrounded assertions of ``high fidelity inferred from $p > 0.05$''. In their place, we provide rigorous subject-level concordance statistics: Lin's Concordance Correlation Coefficient (CCC), Intraclass Correlation Coefficient (ICC), Bland--Altman limits of agreement, and critical amyloid-positivity ($A\beta+$ vs.\ $A\beta-$) classification metrics (Sensitivity, Specificity, Cohen's $\kappa$).

  \item \textbf{Standardized Quantitative Terminology: SUVR vs. SUV (R3-5, R2-9):}
  We standardized all regional uptake terminology to Standardized Uptake Value Ratio (\textbf{SUVR}), clarified reference region normalization (whole cerebellum for AV45, pons/cerebellar gray for FDG), documented the FDR family definitions ($q$-values), and reported Cohen's $d$ effect sizes for all regional comparisons.

  \item \textbf{Strict Subject-Level Isolation \& 2-Fold OOF / 3-Fold Architecture (R3-6, R2-3, R2-6, R2-19):}
  We clarified the two-stage protocol: the generative synthesis model is evaluated via 2-fold Out-Of-Fold (OOF) cross-validation, and downstream classification is evaluated on 3-fold cross-validation with absolute subject-level separation (zero visit leakage). Synthetic PET images evaluated by downstream classifiers were generated exclusively by models that never saw those subjects during training.

  \item \textbf{Systematic Multi-Baseline Comparison \& Ablations (R3-7, R3-9, R3-10, R3-12, R2-17):}
  We added multi-model baselines including tabular MLP, direct regression U-Net, and SwinUNETR baselines under strictly identical splits and preprocessing. All classification metrics are reported with fold-wise mean $\pm$ SD and 95\% confidence intervals.
\end{enumerate}

All modifications are highlighted in \textcolor{red}{red text} in the revised manuscript. Below, we provide our detailed point-by-point responses to each comment raised by the Editor, Reviewer~2, and Reviewer~3.
